% ============================================================================
% Chapter 11 — Animation
% ============================================================================
\chapter{Animation}

\textbf{Headers:} \texttt{animation/AnimationComponents.hpp},
                  \texttt{animation/AnimationSystem.hpp}\\
\textbf{Namespace:} \texttt{Caffeine::Animation}

The animation module provides frame-based sprite animation driven by a
finite-state machine. States hold \texttt{AnimationClip} references and are
connected by \texttt{AnimationTransition} objects whose conditions are
evaluated against named parameters — a design inspired by Unity's and
Unreal's Animator/AnimBlueprint systems. At runtime, the engine's
\texttt{AnimationSystem} evaluates the state machine every frame and updates
the sprite component automatically.

% ─────────────────────────────────────────────────────────────────────────────
\section{AnimationClip}

\textbf{Header:} \texttt{animation/AnimationComponents.hpp}

An \texttt{AnimationClip} is a named sequence of sprite frames.

\begin{lstlisting}
struct FrameRect {
    i32 x, y, w, h;   // pixel rectangle on the atlas texture
};

struct AnimationClip {
    FixedString<32>       name;
    u32                   fps;
    std::vector<FrameRect> frames;
    bool                  loop;

    f32 duration() const;   // computed: frames.size() / (f32)fps
};
\end{lstlisting}

\begin{longtable}{llp{8cm}}
\toprule
\textbf{Field} & \textbf{Type} & \textbf{Description} \\
\midrule
\texttt{name}     & \texttt{FixedString<32>}          & Unique identifier used to reference the clip. \\
\texttt{fps}      & \texttt{u32}                       & Playback frame rate. \\
\texttt{frames}   & \texttt{std::vector<FrameRect>}    & Ordered list of source rectangles on the texture atlas. \\
\texttt{loop}     & \texttt{bool}                      & Whether the clip loops when it reaches the last frame. \\
\texttt{duration()} & \texttt{f32}                     & Returns \texttt{frames.size() / (f32)fps} in seconds. \\
\bottomrule
\end{longtable}

\subsection{Creating a clip}

\begin{lstlisting}
using namespace Caffeine::Animation;

AnimationClip idle;
idle.name = "idle";
idle.fps  = 12;
idle.loop = true;
idle.frames = {
    {  0, 0, 64, 64 },
    { 64, 0, 64, 64 },
    {128, 0, 64, 64 },
    {192, 0, 64, 64 },
};
\end{lstlisting}

% ─────────────────────────────────────────────────────────────────────────────
\section{Parameter System}

Transitions between states are driven by \emph{named parameters} stored in
the \texttt{Animator} component. Setting a parameter at runtime causes the
state machine to evaluate its outgoing transitions on the next update.

\subsection{ParameterType}

\begin{lstlisting}
enum class ParameterType : u8 {
    Bool,
    Float,
    Int,
    Trigger,
};
\end{lstlisting}

A \texttt{Trigger} behaves like a boolean that is automatically reset to
\texttt{false} after it causes a transition to fire.

\subsection{AnimatorParameter}

\begin{lstlisting}
struct AnimatorParameter {
    FixedString<32> name;
    ParameterType   type;

    bool  boolValue;
    f32   floatValue;
    i32   intValue;
    bool  triggered;    // used internally for Trigger type
};
\end{lstlisting}

\subsection{ConditionOperator}

\begin{lstlisting}
enum class ConditionOperator : u8 {
    Equals,
    NotEquals,
    Greater,
    Less,
    GreaterOrEqual,
    LessOrEqual,
};
\end{lstlisting}

\texttt{Equals} / \texttt{NotEquals} work with Bool, Int, and Trigger
parameters. The comparison operators (\texttt{Greater}, \texttt{Less}, etc.)
work with Float and Int parameters.

\subsection{TransitionCondition}

\begin{lstlisting}
struct TransitionCondition {
    FixedString<32>    parameterName;
    ConditionOperator  op;

    bool  boolValue;
    f32   floatValue;
    i32   intValue;
};
\end{lstlisting}

Each condition names a parameter in the \texttt{Animator} and specifies the
comparison that must be true for the transition to be allowed. All conditions
in a transition must be satisfied simultaneously (logical \texttt{AND}).

% ─────────────────────────────────────────────────────────────────────────────
\section{AnimationState and AnimationTransition}

\subsection{AnimationTransition}

\begin{lstlisting}
struct AnimationTransition {
    FixedString<32>                   toState;
    std::vector<TransitionCondition>  conditions;
    f32                               blendTime;
    bool                              hasExitTime;

    // Legacy: function-based condition (still evaluated when conditions is empty)
    std::function<bool()>             legacyCondition;
};
\end{lstlisting}

\begin{longtable}{llp{7.5cm}}
\toprule
\textbf{Field}     & \textbf{Type}   & \textbf{Description} \\
\midrule
\texttt{toState}        & \texttt{FixedString<32>}                  & Name of the destination state. \\
\texttt{conditions}     & \texttt{std::vector<TransitionCondition>} & All conditions must pass for the transition to fire. \\
\texttt{blendTime}      & \texttt{f32}                              & Cross-fade duration in seconds. \\
\texttt{hasExitTime}    & \texttt{bool}                             & When \texttt{true}, the transition can fire only after the current clip finishes. \\
\texttt{legacyCondition}& \texttt{std::function<bool()>}            & Fallback used when \texttt{conditions} is empty. \\
\bottomrule
\end{longtable}

\subsection{AnimationState}

\begin{lstlisting}
struct AnimationState {
    FixedString<32>                    name;
    const AnimationClip*               clip;
    f32                                speed;
    std::vector<AnimationTransition>   transitions;
};
\end{lstlisting}

States are identified by name and reference a clip. Multiple outgoing
transitions can be defined; they are evaluated in order and the first
satisfied transition fires.

% ─────────────────────────────────────────────────────────────────────────────
\section{Animator Component}

\texttt{Animator} is an ECS component that holds the complete state machine
for an entity.

\begin{lstlisting}
struct Animator {
    HashMap<FixedString<32>, AnimationState> states;

    FixedString<32> currentState;
    FixedString<32> previousState;

    f32  timeInState;
    f32  blendWeight;
    f32  playbackScale;
    bool paused;

    std::vector<AnimatorParameter> parameters;

    // Parameter helpers (inline convenience wrappers over the vector above)
    void addParameter(const char* name, ParameterType type);

    void setBool (const char* name, bool  value);
    void setFloat(const char* name, f32   value);
    void setInt  (const char* name, i32   value);
    void setTrigger(const char* name);

    bool getBool (const char* name) const;
    f32  getFloat(const char* name) const;
    i32  getInt  (const char* name) const;
};
\end{lstlisting}

\begin{longtable}{llp{7.5cm}}
\toprule
\textbf{Field}         & \textbf{Type}   & \textbf{Description} \\
\midrule
\texttt{states}        & \texttt{HashMap} & All states keyed by name. \\
\texttt{currentState}  & \texttt{FixedString<32>} & Name of the active state. \\
\texttt{previousState} & \texttt{FixedString<32>} & Name of the state that was active before the last transition. \\
\texttt{timeInState}   & \texttt{f32}    & Time in seconds the entity has spent in \texttt{currentState}. \\
\texttt{blendWeight}   & \texttt{f32}    & Cross-fade weight (\texttt{0.0}–\texttt{1.0}) used during transitions. \\
\texttt{playbackScale} & \texttt{f32}    & Global speed multiplier applied on top of each state's \texttt{speed}. \\
\texttt{paused}        & \texttt{bool}   & When \texttt{true}, the state machine is frozen. \\
\texttt{parameters}    & \texttt{vector} & Named parameters evaluated by transitions. \\
\bottomrule
\end{longtable}

% ─────────────────────────────────────────────────────────────────────────────
\section{AnimationSystem}

\textbf{Header:} \texttt{animation/AnimationSystem.hpp}\\
\textbf{Base class:} \texttt{ECS::ISystem}

\texttt{AnimationSystem} processes every entity that has both an
\texttt{Animator} and a sprite component. It advances \texttt{timeInState},
evaluates all outgoing transitions, fires the first satisfied one, and updates
the sprite's source rectangle to reflect the current frame.

\begin{lstlisting}
class AnimationSystem : public ECS::ISystem {
public:
    // Called automatically by the engine every frame.
    void onUpdate(ECS::World& world, f32 dt) override;

    // ── Playback control ──────────────────────────────────────────────────
    void play   (ECS::World& world, ECS::Entity e, const char* stateName);
    void pause  (ECS::World& world, ECS::Entity e);
    void resume (ECS::World& world, ECS::Entity e);
    void setSpeed(ECS::World& world, ECS::Entity e, f32 speed);
    bool isPlaying(ECS::World& world, ECS::Entity e,
                   const char* stateName) const;

    // ── Parameter write ───────────────────────────────────────────────────
    void addParameter(ECS::World& world, ECS::Entity e,
                      const char* name, ParameterType type);

    void setBool   (ECS::World& world, ECS::Entity e, const char* name, bool  v);
    void setFloat  (ECS::World& world, ECS::Entity e, const char* name, f32   v);
    void setInt    (ECS::World& world, ECS::Entity e, const char* name, i32   v);
    void setTrigger(ECS::World& world, ECS::Entity e, const char* name);

    // ── Parameter read ────────────────────────────────────────────────────
    bool getBool (ECS::World& world, ECS::Entity e, const char* name) const;
    f32  getFloat(ECS::World& world, ECS::Entity e, const char* name) const;
    i32  getInt  (ECS::World& world, ECS::Entity e, const char* name) const;
};
\end{lstlisting}

\subsection{Playback methods}

\begin{longtable}{lp{10cm}}
\toprule
\textbf{Method} & \textbf{Description} \\
\midrule
\texttt{play(world, e, name)}      & Immediately transitions to the named state. Resets \texttt{timeInState}. \\
\texttt{pause(world, e)}           & Freezes playback. Sets \texttt{Animator::paused = true}. \\
\texttt{resume(world, e)}          & Resumes playback from where it was paused. \\
\texttt{setSpeed(world, e, speed)} & Sets \texttt{playbackScale}. Values above 1 speed up, below 1 slow down. \\
\texttt{isPlaying(world, e, name)} & Returns \texttt{true} if the entity is currently in the named state and not paused. \\
\bottomrule
\end{longtable}

\subsection{Parameter methods}

\begin{longtable}{lp{9.5cm}}
\toprule
\textbf{Method} & \textbf{Description} \\
\midrule
\texttt{addParameter(world, e, name, type)} & Registers a new named parameter. Call once during setup. \\
\texttt{setBool(world, e, name, v)}         & Sets a Bool parameter. \\
\texttt{setFloat(world, e, name, v)}        & Sets a Float parameter. \\
\texttt{setInt(world, e, name, v)}          & Sets an Int parameter. \\
\texttt{setTrigger(world, e, name)}         & Sets a Trigger; it is cleared automatically after the transition fires. \\
\texttt{getBool(world, e, name)}            & Returns the current value of a Bool parameter. \\
\texttt{getFloat(world, e, name)}           & Returns the current value of a Float parameter. \\
\texttt{getInt(world, e, name)}             & Returns the current value of an Int parameter. \\
\bottomrule
\end{longtable}

% ─────────────────────────────────────────────────────────────────────────────
\section{Setting Up a Character Animator}

The following example constructs a complete two-state machine (idle / run)
for a player entity, with a Float parameter controlling the transition.

\begin{lstlisting}
#include "animation/AnimationComponents.hpp"
#include "animation/AnimationSystem.hpp"
#include "ecs/World.hpp"

using namespace Caffeine::Animation;

void setupPlayerAnimator(ECS::World& world, ECS::Entity player,
                         AnimationSystem& animSys,
                         const AnimationClip& idleClip,
                         const AnimationClip& runClip)
{
    // 1. Build the state machine
    AnimationState idle;
    idle.name  = "idle";
    idle.clip  = &idleClip;
    idle.speed = 1.0f;

    AnimationState run;
    run.name  = "run";
    run.clip  = &runClip;
    run.speed = 1.0f;

    // 2. Add transitions
    {
        TransitionCondition toRun;
        toRun.parameterName = "speed";
        toRun.op            = ConditionOperator::Greater;
        toRun.floatValue    = 0.1f;

        AnimationTransition idleToRun;
        idleToRun.toState   = "run";
        idleToRun.blendTime = 0.1f;
        idleToRun.conditions.push_back(toRun);
        idle.transitions.push_back(idleToRun);
    }
    {
        TransitionCondition toIdle;
        toIdle.parameterName = "speed";
        toIdle.op            = ConditionOperator::LessOrEqual;
        toIdle.floatValue    = 0.1f;

        AnimationTransition runToIdle;
        runToIdle.toState   = "idle";
        runToIdle.blendTime = 0.1f;
        runToIdle.conditions.push_back(toIdle);
        run.transitions.push_back(runToIdle);
    }

    // 3. Attach the Animator component
    auto& anim = world.add<Animator>(player);
    anim.states["idle"] = idle;
    anim.states["run"]  = run;

    // 4. Register parameters
    animSys.addParameter(world, player, "speed", ParameterType::Float);

    // 5. Start in idle state
    animSys.play(world, player, "idle");
}
\end{lstlisting}

\subsection{Updating the parameter in a script}

\begin{lstlisting}
void PlayerController::onUpdate(ECS::Entity entity,
                                ECS::World& world, f32 dt)
{
    f32 vel = computeHorizontalVelocity();   // game-specific
    animSys.setFloat(world, entity, "speed", std::abs(vel));
}
\end{lstlisting}

The \texttt{AnimationSystem} will pick up the change on the next
\texttt{onUpdate} call and fire the appropriate transition automatically.

% ─────────────────────────────────────────────────────────────────────────────
\section{Using Triggers for One-Shot Animations}

Triggers are ideal for actions that happen once (jump, attack, hit).

\begin{lstlisting}
// Setup: add a Trigger parameter and connect a transition
animSys.addParameter(world, player, "attack", ParameterType::Trigger);

TransitionCondition cond;
cond.parameterName = "attack";
cond.op            = ConditionOperator::Equals;
cond.boolValue     = true;   // Trigger reads as bool internally

AnimationTransition idleToAttack;
idleToAttack.toState      = "attack";
idleToAttack.blendTime    = 0.05f;
idleToAttack.hasExitTime  = false;
idleToAttack.conditions.push_back(cond);
idle.transitions.push_back(idleToAttack);

// At runtime: fire the attack animation once
animSys.setTrigger(world, player, "attack");
// The trigger is cleared automatically after the transition fires.
\end{lstlisting}

% ─────────────────────────────────────────────────────────────────────────────
\section{Editor Windows}

\subsection{Animation Timeline}

The \textbf{Animation Timeline} window (\texttt{editor/AnimationTimeline.hpp})
provides a visual editor for individual \texttt{AnimationClip} objects.

\begin{itemize}[noitemsep]
  \item \textbf{Ruler} — Displays time tick marks derived from the clip's
        duration. Click anywhere on the ruler to seek the playhead to that
        time.
  \item \textbf{Playhead} — A red vertical line with a triangle marker on the
        ruler indicates \texttt{currentTime}. Drag it to scrub the clip.
  \item \textbf{Tracks} — Each track corresponds to a channel (e.g., frame
        index, a property). Rows alternate in shade for readability; the
        track name is rendered as a coloured label on the left.
  \item \textbf{Keyframes} — Rendered as filled diamonds. Left-click and drag
        to reposition. Right-click to delete.
  \item \textbf{Buttons} — \texttt{+ Track} adds a new channel,
        \texttt{+ Key} inserts a keyframe at the current playhead position,
        \texttt{Del Key} removes the selected keyframe.
\end{itemize}

\subsection{Animator Controller}

The \textbf{Animator Controller} window
(\texttt{editor/AnimatorController.hpp}) provides a visual state-machine
editor comparable to Unity's Animator window.

\begin{itemize}[noitemsep]
  \item \textbf{Canvas} — An infinite scrollable grid. Pan by dragging with
        the middle mouse button.
  \item \textbf{State nodes} — Draggable boxes. The default state is
        highlighted in a distinct colour; the currently active state is shown
        in a third colour.
  \item \textbf{Transitions} — Drawn as directed lines with a triangular
        arrowhead at the midpoint indicating direction.
  \item \textbf{Right-click canvas} — Opens an \emph{Add State} modal to
        create a new state node.
  \item \textbf{Right-click node} — Context menu with options:
        \emph{Set as Default}, \emph{Add Transition}, \emph{Delete State}.
  \item \textbf{Parameters panel} (left) — Lists all \texttt{AnimatorParameter}
        entries; values can be edited live to test transitions at runtime.
  \item \textbf{Inspector panel} (right) — Shows the selected state's
        \texttt{speed}, \texttt{clip}, outgoing transitions, and each
        transition's conditions.
  \item \texttt{setAnimator(animator*)} — Bind a live \texttt{Animator}
        component to the window; state nodes are auto-positioned on a grid.
\end{itemize}

% ─────────────────────────────────────────────────────────────────────────────
\section{Best Practices}

\begin{itemize}[noitemsep]
  \item Name parameters and states consistently — use \texttt{camelCase}
        for parameters (\texttt{"isGrounded"}, \texttt{"speed"}) and
        \texttt{PascalCase} for state names (\texttt{"Idle"}, \texttt{"Run"}).
  \item Prefer Float/Bool parameters over Triggers for conditions that reflect
        ongoing game state. Triggers are appropriate only for discrete,
        one-shot events.
  \item Keep \texttt{blendTime} short (0.05–0.15 s) for responsive
        characters and longer (0.2–0.4 s) for cinematic transitions.
  \item Set \texttt{hasExitTime = true} only for transitions that must wait
        for the current clip to finish, such as an attack returning to idle.
  \item Call \texttt{addParameter} once during entity setup, not every frame.
        Reading or writing a non-existent parameter is a no-op, but registering
        it late may miss the first frame's transitions.
\end{itemize}
